\documentclass[11pt]{amsart}
\usepackage[T1]{fontenc}
\usepackage[utf8]{inputenc}
\usepackage{amssymb}
\usepackage{array}
\usepackage[margin=1.1in]{geometry}
\usepackage[hidelinks]{hyperref}
\newcount\sthh\newcount\stmm
\sthh=\time \divide\sthh by 60 \stmm=-\sthh \multiply\stmm by 60 \advance\stmm by \time
\newcommand{\stamptime}{\the\sthh:\ifnum\stmm<10 0\fi\the\stmm}
\AddToHook{shipout/foreground}{\put(0.5\paperwidth,-\paperheight+0.55cm){\makebox[0pt]{\tiny\ttfamily unit/doc/\jobname.tex \ --- \ compiled \today, \stamptime}}}
\newtheorem{theorem}{Theorem}
\newtheorem{proposition}[theorem]{Proposition}
\newtheorem{lemma}[theorem]{Lemma}
\newtheorem{corollary}[theorem]{Corollary}
\theoremstyle{definition}
\newtheorem{definition}[theorem]{Definition}
\newtheorem{statement}[theorem]{Statement}
\theoremstyle{remark}
\newtheorem{remark}[theorem]{Remark}
\newcommand{\Om}{\Omega}
\newcommand{\Z}{\mathbb{Z}}
\newcommand{\code}[1]{\texttt{#1}}
\newcommand{\Map}{\mathrm{Map}_*}
\newcommand{\hc}[2]{[#1,#2]_*}
\newcommand{\Fib}{\mathrm{Fib}}
\newcommand{\fpr}{\mathsf{fpr}}
\newcommand{\tr}[1]{\|#1\|}
\newcommand{\cls}[1]{[#1]}
\newcommand{\kb}{\bar\kappa}
\newcommand{\lp}{\mathsf{loop}}
\newcommand{\Unl}{\mathsf{Unloop}}
\newcommand{\SCt}{\mathsf{SC}_3}
\newcommand{\SCf}{\mathsf{SC}_4}
\newcommand{\Tri}{\mathsf{Triples}}
\newcommand{\pteq}{\simeq_*}
\newcommand{\refl}{\mathsf{refl}}
\title{$\pi_4(S^3)=\Z/2$ through the delooping embedding:\\ the main ideas}
\date{28 October 2026}
\begin{document}
\maketitle
\sloppy

\noindent\textbf{Attribution.} The ingredients are David W\"arn's: a
candidate Postnikov stage $BR$ identified by a \emph{marked universal property},
and the fact that \emph{delooping is an embedding} of pointed $2$-connected
$4$-types into pointed simply connected $3$-types. The proof below uses them,
but it is \emph{not} the route W\"arn had in mind, which goes through a
classification of arbitrary maps between triples; we bypass that step with
specialised mapping arguments (Section~\ref{s:route}). The delooping machinery (Eilenberg--MacLane spaces
as iterated torsors, the criterion for a map of loop spaces to deloop) is
W\"arn's~\cite{DW1}. The classification of pointed simply connected $3$-types by
triples is Hart--Milner's (Owen Milner's equivalence)~\cite{HM}; the
development re-proves it in that case. The cup-square
calculations used underneath (cross term, bilinear classification, Lemma~N)
are those of the earlier proof, \code{doc/pi4-s3.pdf}.

\smallskip\noindent\textbf{The result.} In univalent type theory with
replacement, assuming only the universal property of $S^3$,
\[\pi_4(S^3)\cong\Z/2\quad\text{as groups}\]
(\code{Pi4Direct.Pi4S3.pi4-S3}). It is a second proof of the theorem of
\code{doc/pi4-s3.pdf} (\code{Pi4Sphere.Pi4S3.pi4-S3}): the two share the
candidate $BR$, the marked class $s$, and the final Yoneda argument, but differ
in how the universal property at $\tr{S^3}_4$ is obtained. The earlier proof
builds a Postnikov presentation of $\tr{S^3}_4$; the present one proves the
universal property at \emph{every} pointed $2$-connected $4$-type, from a
restricted $2$-torsion lemma, and needs no presentation.

\section{Where the argument goes}

\noindent\textbf{Notation.} $\Map(X,Y)$ is the type of pointed maps,
$\hc XY$ its set truncation, $c;d$ composition of classes. $K(A,n)$ are the
Eilenberg--MacLane spaces, $e_m:\Om K(A,m{+}1)\pteq K(A,m)$. $\kappa:K(\Z,2)\to_*
K(\Z,4)$ is the constructed cup square, $\rho:\Z\to\Z/2$ the parity, and
$\kb=\kappa;K(\rho,4)$. For a pointed map $k$, $\Fib k$ is its fibre and
$\fpr:\Fib k\to_*$ (source of $k$) the projection.

\smallskip\noindent\textbf{The candidate and the final step.} Let $R=\Fib\kb$, a
pointed simply connected $3$-type with $\pi_3R=\Z/2$. The cross term of
$\kb$ vanishes, so $\kb$ deloops to $\hat F:K(\Z,3)\to_*K(\Z/2,5)$ and
$BR=\Fib\hat F$ is a pointed $2$-connected $4$-type with $\Om BR\pteq R$; in
particular $\pi_4BR\cong\pi_3R\cong\Z/2$. There is a class $s\in\hc{S^3}{BR}$
(\code{SourceS3.Final.s}). Write
\[
  \mathsf M(Y):\qquad s^*:\hc{BR}{Y}\to\hc{S^3}{Y},\ c\mapsto s;c,\quad\text{is a bijection.}
\]
If $\mathsf M(BR)$ and $\mathsf M(T)$ hold for $T=\tr{S^3}_4$, a two-object
Yoneda argument on classes gives inverse classes $T\leftrightarrows BR$, hence
$\pi_4S^3\cong\pi_4T\cong\pi_4BR\cong\Z/2$ (\code{Pi4Sphere.Two4}).
$\mathsf M(BR)$ is the mapping theorem for targets $\Fib K$ of level-$3$
triples (\code{SourceS3.Final.rep-fib}). The whole difficulty is $\mathsf M(T)$.

\smallskip\noindent\textbf{The new route.} We prove $\mathsf M(Y)$ for every
pointed $2$-connected $4$-type $Y$:
\begin{enumerate}
\item every such $Y$ has $\Om Y\pteq\Fib k$ for a \emph{deloopable} triple
$(A,B,k)$ (Section~\ref{s:emb});
\item $\mathsf M(Y)$ is equivalent to its looped form, a statement about maps
into $\Om Y$ (Section~\ref{s:loop});
\item the looped form reduces, by a lifting lemma along $\fpr$, to one
numerical fact about deloopable triples (Section~\ref{s:red});
\item that fact is Statement~T, the restricted $2$-torsion lemma
(Section~\ref{s:T}): a cup-square calculation in which deloopability enters only
through the multiplicativity of one loop map.
\end{enumerate}

\section{Delooping is an embedding}\label{s:emb}

Let $\SCf$ be the type of pointed $2$-connected $4$-types and $\SCt$ that of
pointed simply connected $3$-types; looping is a map $\Om:\SCf\to\SCt$.

\begin{lemma}[looping of maps]\label{lem:loop}
If $X$ is $2$-connected and $Y$ is $4$-truncated, then
$\lp:\Map(X,Y)\to\Map(\Om X,\Om Y)$ is an equivalence.
\end{lemma}

This is W\"arn's delooping criterion~\cite{DW1}: the type of deloopings of a
map of loop spaces is contractible under these connectivity and truncation
bounds (\code{SourceS3.LoopEquiv}).

\begin{theorem}[W\"arn]\label{thm:emb}
$\Om:\SCf\to\SCt$ is an embedding: its fibres are propositions.
\end{theorem}

\noindent By Lemma~\ref{lem:loop} and Whitehead, looping of pointed
equivalences $(u\pteq v)\to(\Om u\pteq\Om v)$ has a section; so the type of
$v$ with $\Om u\pteq\Om v$ is a retract of the (contractible) type of $v$ with
$u\pteq v$, and the fundamental theorem of identity types gives the claim
(\code{DeloopEmbedding.Emb}).

By the classification of Hart--Milner, $\SCt\simeq\Tri$, where a triple
$(A,B,k)$ consists of abelian groups $A,B$ and $k:K(A,2)\to_*K(B,4)$, realised
as $\Fib k$ (\code{PostnikovClassification.classification}). Call a triple
\emph{deloopable} if $\Fib k$ is in the image of $\Om$; by
Theorem~\ref{thm:emb} this is a proposition, and every $Y\in\SCf$ is a delooping
of the triple classifying $\Om Y$:
\[
  \textstyle\sum_{t:\Tri}\mathsf{Deloopable}(t)\ \simeq\ \SCf .
\]
So it suffices to understand maps into $\Om Y\pteq\Fib k$ for deloopable
triples. Note that being deloopable is weaker, a priori, than $k$ itself being
the looping of a map $K(A,3)\to K(B,5)$; we never need to decide this.

\section{The looped form of M}\label{s:loop}

Let $r=(\lp s);\cls{\Om BR\pteq R}\in\hc{\Om S^3}{R}$, the looping of the marked
class, viewed with target $R$ (\code{MarkLoop.LMark.r}).

\begin{proposition}[L]\label{prop:L}
For $Y\in\SCf$: $\mathsf M(Y)$ holds iff
$r^*:\hc{R}{\Om Y}\to\hc{\Om S^3}{\Om Y}$ is a bijection.
\end{proposition}

\noindent Looping commutes with precomposition, and by
Lemma~\ref{lem:loop} looping is a bijection
$\hc{BR}{Y}\to\hc{\Om BR}{\Om Y}$ and $\hc{S^3}{Y}\to\hc{\Om S^3}{\Om Y}$
($BR$ and $S^3$ are $2$-connected). Two applications of 2-out-of-3 conclude
(\code{MarkLoop.LMark.At.to-loop}, \code{from-loop}).

The point of the looped form is that its target $\Om Y\pteq\Fib k$ is a fibre,
and maps into a fibre are analysed by a lifting lemma.

\section{Reduction to a numerical fact}\label{s:red}

Fix a deloopable triple $(A,B,k)$. Postcomposition with $\fpr$ gives, for every
pointed source $Q$, a map $q_Q:\hc{Q}{\Fib k}\to\hc{Q}{K(A,2)}$, and the square
\[
  q_{\Om S^3}(r;c)=r;q_R(c)
\]
commutes. We show that $q_R$, $q_{\Om S^3}$ and $r^*:\hc R{K(A,2)}\to\hc{\Om
S^3}{K(A,2)}$ are bijections; then $r^*$ is a bijection into $\Fib k$ by
2-out-of-3, and into $\Om Y$ by transport along $\Om Y\pteq\Fib k$; Proposition~L
gives $\mathsf M(Y)$.

\smallskip\noindent\textbf{$r^*$ into $K(A,2)$.} $K(A,3)\in\SCf$ and
$\mathsf M(K(A,3))$ holds (it is part of the construction of $s$:
$s;\cls p=\cls{\sigma_3}$ with $p:BR\to K(\Z,3)$ and $\sigma_3$ the generator).
Proposition~L turns it into the bijectivity of $r^*$ into $\Om K(A,3)\pteq
K(A,2)$.

\smallskip\noindent\textbf{The lifting lemma.} For a source $Q$, $q_Q$ is a
bijection provided (E) for every $F:Q\to_*K(A,2)$, $F;k$ is merely null, and
(U) every map $Q\to_*\Om K(B,4)$ is merely null (\code{MapFib.Lifting}). For
$Q=\Om S^3$, (U) holds because, by Lemma~\ref{lem:loop}, such a map is the
looping of a map $S^3\to_*K(B,4)$, which is null since $\Om^3K(B,4)$ is
connected; and (E) follows from (E) at $R$, because every class from $\Om S^3$
is $r;c$ with $c$ from $R$ (bijectivity of $r^*$ into $K(A,2)$). For $Q=R$,
(U) is $\hc{R}{K(B,3)}=0$ (proved for the earlier proof, from (H3) and (H4)).
Everything therefore rests on:

\smallskip\noindent\textbf{(E) at $R$.} \emph{For every $F:R\to_*K(A,2)$, $F;k$
is merely null.}

\smallskip\noindent The fibres of $p_R=\fpr:R\to K(\Z,2)$ are $2$-connected,
so $F$ factors as $p_R;G$ for some $G:K(\Z,2)\to_*K(A,2)$. By (H4),
$\cls{G;k}=\Phi_B(h)=\cls{\kappa;K(h,4)}$ for a unique $h:\Z\to B$. If
$2h(1)=0$, then $h=\rho;g$ for some $g:\Z/2\to B$, so
$\cls{G;k}=\cls{\kb;K(g,4)}$, and $F;k=p_R;\kb;K(g,4)$ is merely null because
$p_R;\kb$ is. So (E) at $R$ follows from:

\begin{statement}[T]\label{st:T}
Let $(A,B,k)$ be deloopable, $G:K(\Z,2)\to_*K(A,2)$, and $\cls{G;k}=\Phi_B(h)$.
Then $2h(1)=0$.
\end{statement}

\noindent This is where deloopability is used, and the only place.

\section{Statement T}\label{s:T}

\noindent\textbf{The cross term.} For $f:\Om K(\Z,3)\to_*\Om K(B,5)$, the
\emph{cross term} measures the failure of $f$ to be a homomorphism:
\[
  \Delta^f_q(p)=(fp)^{-1}\cdot f(p\cdot q)\cdot(fq)^{-1}.
\]
Packaged through $e_2$ and $e_4$, it is a bipointed map
$\delta_f:K(\Z,2)\to_*\Map(K(\Z,2),K(B,4))$ (\code{CrossTerm.CrossGen}). Two
facts from the earlier proof are used as black boxes:
\begin{itemize}
\item[(a)] if $\delta_f=0$ then $f$ deloops: $f=\lp\hat f$ (the cells argument,
\code{DeloopR.UnloopGen});
\item[(b)] Lemma~N: if $f_F=e_2;F;e_4^{-1}$ deloops, where $F=\kappa;K(h,4)$,
then $2h(1)=0$ (\code{Necessity.Conv}).
\end{itemize}
So Statement~T follows once we show $\delta_{f'}=0$ for $f'=f_{G;k}$; then
$f'$ deloops, $G;k=\kappa;K(h,4)$ merely, and (b) applies. In fact we obtain
the stronger conclusion that $f_{G;k}$ deloops.

\smallskip\noindent\textbf{Detecting the cross term.} By the bilinear
classification (Theorem~B of \code{doc/h4.pdf}), a bipointed map $c$ is
determined by its \emph{evaluation on the generator}: $E(c)(x)\in\Om^2K(B,4)$,
the looping of $c$ evaluated at the generator of the circle. Here this takes a
concrete form: with $\hat\gamma_0\in\Om^3K(\Z,3)$ the generator and
$\lambda_\gamma(m)=(\lp^2m)(\gamma)$,
\[
  E(\delta_f)(x)=(\lp^2e_4)\big(\lambda_{\hat\gamma_0}(\Delta^f_q)\big),\qquad q=e_2^{-1}x .
\]
So $\delta_f=0$ as soon as $\lambda_{\hat\gamma_0}(\Delta^f_q)=\refl$ for all
$q$. Only the underlying function of $\Delta^f_q$ matters: $\lambda$ is
insensitive to the pointing (an Eckmann--Hilton argument at two levels).

\smallskip\noindent\textbf{The H-space argument.} Let $\epsilon:\Om Y\pteq\Fib
k$ witness deloopability and
\[
  \pi=\epsilon;\fpr;e_2^{-1}:\Om Y\to_*\Om K(A,3),\qquad f_k=e_2;k;e_4^{-1}.
\]
By Lemma~\ref{lem:loop}, $\pi=\lp P$ for some $P:Y\to K(A,3)$; hence $\pi$ is
\emph{multiplicative}. This is the one consequence of deloopability we use.
Three observations:
\begin{enumerate}
\item $\pi;f_k=\epsilon;(\fpr;k);e_4^{-1}$ is merely null, since $\fpr;k$ is.
\item For a multiplicative $L$, $\Delta^{L;f}_{q}(p)=\Delta^f_{Lq}(Lp)$. Hence
$\Delta^{f_k}_{\pi b}\circ\pi=\Delta^{\pi;f_k}_b$, which is (pointwise) $\Delta^0_b=\refl$.
\item $f'=\hat G;f_k$ with $\hat G=e_2;G;e_2^{-1}=\lp K(a,3)$ (since
$G=K(a,2)$), again multiplicative.
\end{enumerate}
Choose $\gamma'\in\Om^3Y$ with $(\lp^2\pi)\gamma'=(\lp^2\hat G)\hat\gamma_0$; it
exists merely because $\lp^2\fpr$ is surjective (the fibre of $\fpr$ is
$2$-connected). Define $\varphi(z)=\lambda_{\gamma'}(\pi;\Delta^{f_k}_z)$ for
$z\in\Om K(A,3)$. By (1) and (2), $\varphi$ vanishes on the image of $\pi$, i.e.\
along $\fpr$; since its target $\Om^3K(B,5)\simeq K(B,2)$ is a $2$-type and the
fibres of $\fpr$ are $2$-connected, $\varphi$ vanishes everywhere. Then
\[
  \lambda_{\hat\gamma_0}(\Delta^{f'}_q)
  =\lambda_{\hat\gamma_0}(\hat G;\Delta^{f_k}_{\hat Gq})
  =\lambda_{(\lp^2\hat G)\hat\gamma_0}(\Delta^{f_k}_{\hat Gq})
  =\lambda_{\gamma'}(\pi;\Delta^{f_k}_{\hat Gq})=\varphi(\hat Gq)=\refl ,
\]
using (3), naturality of $\lambda$, and the choice of $\gamma'$. Hence
$E(\delta_{f'})=E(0)$ pointwise and $\delta_{f'}=0$.

The two mere choices (the null-homotopy in (1) and $\gamma'$) are eliminated
into the goal $\Theta\,\delta_{f'}=\Theta\,0$, an equality in the \emph{set}
$\Map(K(\Z,2),K(B,2))$ ($\Theta$ is $E$ followed by $\Om^2K(B,4)\simeq K(B,2)$);
pointwise equality of $E$ alone would not be a proposition.

\begin{remark}[classically]
The argument says: the cross term of $k$, a class in
$\tilde H^4(K(A,2)\wedge K(A,2);B)$, vanishes after pullback to $\Fib k\wedge\Fib k$
because $\Fib k$ is an H-space over $K(A,2)$ (it is a loop space, and $\fpr$ is
a loop map up to $e_2$), and pullback is injective on the $\pi_2\otimes\pi_2$
part, which is what $E$ detects. The cross term of $\kappa;K(h,4)$ is
$2h\,\iota\otimes\iota$; Lemma~N is used instead of this computation.
\end{remark}

\section{Assembly}

Putting the pieces together: for every $Y\in\SCf$, $\mathsf M(Y)$ holds
(\code{MarkDirect.Direct.marked}). At $T=\tr{S^3}_4$, which is $2$-connected and
$4$-truncated, this is the universal property needed by the Yoneda argument;
with $\mathsf M(BR)$ and $\pi_4BR\cong\Z/2$ it gives
\[
  \pi_4(S^3)\cong\pi_4\tr{S^3}_4\cong\pi_4(BR)\cong\Z/2 .
\]

\section{Relation to the intended route}\label{s:route}

\noindent\textbf{What is re-proved.} The development re-proves Owen Milner's
classification equivalence in its simply connected $3$-type case,
\[
  \SCt\ \simeq\ \textstyle\sum_{A,B:\mathrm{Ab}}\Map\big(K(A,2),K(B,4)\big),
\]
following Hart--Milner (Lemma~40, Theorem~41, Note~42 at $n=2$, over a
simply connected base) (\code{PostnikovClassification.classification}). The
argument above \emph{uses} this classification, on objects and equivalences
only, together with W\"arn's looping embedding (Theorem~\ref{thm:emb}) and the
direct $2$-torsion argument (Statement~T); it does not re-prove the
classification a second time.

\smallskip\noindent\textbf{What is not done.} The development does not
classify \emph{arbitrary pointed maps} between realised triples, i.e.\ it does
not describe $\Map(\Fib k,\Fib k')$ through maps of coefficient groups
together with compatible homotopies between the $k$-invariants, and so it does
not obtain the resulting equivalence of homotopy categories.

\smallskip\noindent\textbf{Consequence.} The route W\"arn had in mind
includes exactly that step: describe arbitrary maps between realised triples
in terms of the triple data, then identify the universal object of the sphere
on the triple side. The checked proof keeps his classification-and-looping
perspective but bypasses the general step: the maps it needs are handled by
specialised mapping theorems (the lifting lemma along $\fpr$ at the sources
$R$, $\Om S^3$, and, for $\mathsf M(BR)$, $BR$ and $S^3$; Section~\ref{s:red}).
It is therefore a different implementation of the idea, not a formalisation of
the intended proof.

\smallskip\noindent\textbf{What remains for the intended formulation.}
(i)~The general correspondence on morphisms: pointed maps $\Fib k\to_*\Fib k'$
(up to homotopy, or as a type) in terms of group maps and compatible
homotopies of $k$-invariants, compatible with identities and composition, so
that the classification becomes an equivalence of (homotopy) categories, and
its looping counterpart for $\SCf$. (ii)~A reformulation of the
universal-property argument through it: the marked property of $BR$ stated
and proved on the triple side.

\section{What is formalised}\label{s:formal}

Everything above is checked in Unit, development \code{sag/pi3-s2-postnikov}
(\code{./segc sag/pi3-s2-postnikov all}). The new modules are:

\smallskip
\noindent
\begin{tabular}{@{}>{\raggedright\arraybackslash}p{0.2\textwidth}>{\raggedright\arraybackslash}p{0.76\textwidth}@{}}
\code{DeloopEmbedding} & Theorem~\ref{thm:emb}, deloopable triples, $\sum_t\mathsf{Deloopable}(t)\simeq\SCf$\\
\code{MarkLoop} & looping on classes; the class $r$; Proposition~L\\
\code{CrossDeloop} & $\Delta^f_q$, $\lambda$, pointing irrelevance, the formula for $E(\delta_f)$\\
\code{TwoTorsion} & Statement~T: $f_{G;k}$ deloops (\code{TwoTors.unloop}) and $2h(1)=0$ (\code{TwoTors.Two.two-tors}); no sphere axiom\\
\code{MarkDirect} & (E) at $R$, the lifting lemma at $R$ and $\Om S^3$, $\mathsf M(Y)$ for every $Y\in\SCf$\\
\code{Pi4Direct} & the assembly, \code{Pi4S3.pi4-S3 : GIso[pi4 S3 base, ab-grp Z2ab]}
\end{tabular}

\smallskip\noindent The final theorem depends only on the universal property of
$S^3$ (a module parameter) and the Theory parameters (univalence, replacement).
It does not use the Postnikov presentation of $\tr{S^3}_4$
(\code{Presentation4}): a copy of the development with that module deleted
still checks. Detailed informal proofs, declaration by declaration, are in
\code{doc/delooping-direct.pdf}; the earlier route is \code{doc/pi4-s3.pdf}.

\begin{thebibliography}{9}
\bibitem{DW1} D.~W\"arn. \emph{Eilenberg--MacLane spaces and stabilisation in homotopy
  type theory}. arXiv:2301.03685, 2023.
\bibitem{HM} P.~Hart and O.~Milner. \emph{Classifying 2-groups in homotopy type
  theory}. LICS 2026, LIPIcs vol.~380.
\end{thebibliography}
\end{document}
